% Bagian: computability
% Bab: computability-theory
% Subbagian: coding-computations

\documentclass[../../../include/open-logic-section]{subfiles}

\begin{document}

\olfileid[id]{cmp}{thy}{cod}
\olsection{Pengodean Komputasi}

Dalam setiap model komputasi, hal-hal berikut dapat dilakukan:
\begin{enumerate}
\item Mendeskripsikan \emph{definisi} fungsi-fungsi dapat dihitung secara
  sistematis. Misalnya, spesifikasi mesin Turing, definisi rekursif, atau
  program dalam suatu bahasa pemrograman dapat dipandang sebagai pemberi
  definisi tersebut.
\item Mendeskripsikan rekaman lengkap komputasi suatu fungsi yang diberikan
  oleh suatu definisi, untuk masukan tertentu. Misalnya, komputasi mesin Turing
  dapat dideskripsikan oleh urutan konfigurasi (keadaan mesin dan isi pita)
  pada setiap langkah komputasi.
\item Menguji apakah suatu calon rekaman komputasi benar-benar merekam
  komputasi suatu fungsi dapat dihitung dengan definisi tertentu pada masukan
  tertentu (yang padanya fungsi itu terdefinisi, yakni komputasinya berhenti).
\item Mengekstrak nilai fungsi untuk suatu masukan tertentu dari deskripsi
  rekaman lengkap komputasi semacam itu. Misalnya, isi pita pada langkah
  terakhir suatu komputasi mesin Turing yang berhenti merupakan nilai fungsi
  tersebut.
\end{enumerate}

Melalui pengodean, setiap deskripsi fungsi dapat dihitung dapat diberi
\emph{indeks} numerik sedemikian rupa sehingga instruksinya dapat diperoleh
kembali dari indeks itu dengan cara yang dapat dihitung. Demikian pula,
rekaman lengkap suatu komputasi dapat dikodekan dengan satu bilangan. Relasi
aritmetis yang dihasilkan, yakni ``$s$~mengodekan rekaman komputasi fungsi
berindeks~$e$ untuk masukan~$x$,'' dan fungsi ``keluaran urutan komputasi
berkode~$s$,'' kemudian dapat dihitung; bahkan, keduanya bersifat rekursif
primitif.

Fakta mendasar ini sangat kuat dan memungkinkan kita membuktikan sejumlah
hasil yang mencolok dan penting tentang komputabilitas, terlepas dari model
komputasi yang dipilih.

\end{document}
